\section*{Capítulo \ref{ch:b2:climate-variability} --- Variabilidad climática}

\begin{solution}{exo:b2:climate-variability:1}
(a) $29.27 - 26.71 = \qty{2.56}{\celsius}$ en 2026; $28.79 - 26.73 = \qty{2.06}{\celsius}$ en 1997.
(b) Septiembre de 2026 fue a la vez el mes más cálido, por \qty{0.48}{\celsius}, y la mayor anomalía, por
\qty{0.50}{\celsius}. (c) Cada año se compara con un periodo base de 30 años centrado en él, que se desplaza
cada cinco años; la \omterm{def:b2:climate-variability:modes}{climatología} de septiembre de las décadas de 1980 y 1990 difiere ligeramente de la de
las décadas de 2010 y 2020.
\end{solution}

\begin{solution}{exo:b2:climate-variability:2}
(a) JJA: $(0.79 + 0.95 + 1.28)/3 = 1.007$; JAS: $(0.95 + 1.28 + 1.55)/3 = 1.260$; ASO: $(1.28 +
1.55 + 1.68)/3 = 1.503$; SON: $(1.55 + 1.68 + 2.01)/3 = 1.747$; OND: $(1.68 + 2.01 + 2.02)/3 =
1.903$, todo en \unit{\celsius}. (b) Redondeados a un decimal: 1.0, 1.3, 1.5, 1.7, 1.9, los valores de
la tabla. (c) En ASO 2023, con 1.5.
\end{solution}

\begin{solution}{exo:b2:climate-variability:3}
(a) El índice sigue la desviación de la presión en Tahití menos la de Darwin: por debajo de la media
en Tahití y por encima en Darwin, es negativo. (b) El~Niño. (c) Alisios del este más débiles, incluso
vientos del oeste cerca de la superficie en el oeste, y menos lluvia de lo habitual sobre Indonesia, ya que la
convección profunda se ha desplazado hacia el este.
\end{solution}

\begin{solution}{exo:b2:climate-variability:4}
(a) $804.1/880.6 = \qty{91.3}{\percent}$ de la media: entre el 90 y el \qty{110}{\percent}, una
estación normal, aunque en su parte baja. (b) La estación suma $556 + 659 + 427 + 252 =
\qty{1894}{mm}$; julio aporta $659/1894 = \qty{34.8}{\percent}$, más de un tercio.
\end{solution}

\begin{solution}{exo:b2:climate-variability:5}
(a) $\tau = 1025 \times 3992 \times 50/16.1 = \qty{1.271e7}{s} = \qty{147.1}{\day}$ para
\qty{50}{m}, y el doble, \qty{294.2}{\day}, para \qty{100}{m}. (b) $\mathrm{e}^{-30.44/147.1} =
0.813$ y $\mathrm{e}^{-30.44/294.2} = 0.902$. (c) El tiempo de amortiguamiento es proporcional a $h$: una
capa de mezcla menos profunda almacena menos calor para el mismo flujo y olvida antes una anomalía, así que la
memoria del océano, y la predictibilidad que aporta, disminuyen.
\end{solution}

\begin{solution}{exo:b2:climate-variability:6}
(a) Tres meses son 91.3 días: $1.2 \times \mathrm{e}^{-91.3/147.1} = \qty{0.65}{K}$; seis meses
son 182.6 días: $1.2 \times \mathrm{e}^{-182.6/147.1} = \qty{0.35}{K}$. (b) La anomalía decae
despacio, así que es una estimación de la anomalía del mes siguiente mejor que cero. Pero el forzamiento aleatorio añade
nuevas anomalías cada día, desconocidas en el momento de la predicción; a medida que la anomalía inicial decae,
estas ocupan una parte cada vez mayor, y la predicción deriva hacia la \omterm{def:b2:climate-variability:modes}{climatología}, donde ya no
ayuda.
\end{solution}

\begin{solution}{exo:b2:climate-variability:7}
(a) $160 \times 86\,400 = \qty{1.38e7}{J}$, \qty{13.8}{MJ} por metro cuadrado y por día. (b) La
capa contiene $(1000 - 850) \times 100/9.81 = \qty{1529}{kg}$ de aire por metro cuadrado, con una capacidad
calorífica de $1529 \times 1004.5 = \qty{1.536e6}{J.m^{-2}.K^{-1}}$; se calentaría $1.38 \times
10^7/1.536 \times 10^6 = \qty{9.0}{K}$ al día. (c) La capa radia calor en el infrarrojo, hacia el
espacio y hacia el aire de encima; y el aire cálido asciende y se mezcla con la capa superior, o
lo arrastran los vientos, entre ellos la entrada de aire más fresco que la propia baja térmica atrae.
\end{solution}

\begin{solution}{exo:b2:climate-variability:8}
(a) $(-1.93 - 1.11 - 1.98)/3 = -1.673$, aproximadamente $-1.67$. (b) Es el invierno más bajo del registro.
(c) Una NAO fuertemente negativa: una baja de Islandia y un anticiclón de las Azores débiles, un chorro más débil y una trayectoria de las borrascas
más al sur de lo habitual; un invierno frío y seco en el norte de Europa, húmedo en el sur de Europa,
y más irrupciones de aire frío en el este de Estados Unidos.
\end{solution}

\begin{solution}{exo:b2:climate-variability:9}
(a) $9.3 - 0.4 = \qty{8.9}{mm}$ al día, más de 2.5: en el dominio. (b) En el hemisferio sur
el verano es diciembre--febrero: $6.1 - 4.2 = \qty{1.9}{mm}$ al día, menos de 2.5: fuera de él,
aunque allí llueve más que en la mayoría de los lugares del dominio. (c) $1894/122 = \qty{15.5}{mm}$ al
día; Kerala solo quedaría fuera si su intensidad invernal superara $15.5 - 2.5 = \qty{13.0}{mm}$ al día.
(d) Un \omterm{def:b2:climate-variability:monsoon}{monzón} es una inversión estacional: un lugar húmedo todo el año, como la segunda
estación, no tiene \omterm{def:b2:climate-variability:monsoon}{monzón}, y un lugar seco con una estación lluviosa corta e intensa sí lo tiene.
\end{solution}

\begin{solution}{exo:b2:climate-variability:10}
(a) Derivando, $\ddot T = a\dot H = -abT$: la ecuación de un oscilador armónico de frecuencia
angular $\omega = \sqrt{ab}$ y periodo $2\pi/\sqrt{ab}$. (b) $2\pi/\sqrt{ab} = 4$ da $ab =
(2\pi/4)^2 = \qty{2.47}{yr^{-2}}$. (c) Con $T = u\,\mathrm{e}^{-rt}$, $H = v\,\mathrm{e}^{-rt}$,
$\dot T = (\dot u - ru)\,\mathrm{e}^{-rt} = (av - ru)\,\mathrm{e}^{-rt}$, así que $\dot u = av$, y
del mismo modo $\dot v = -bu$: $(u, v)$ oscila sin amortiguamiento, y $(T, H)$ es esa oscilación
multiplicada por $\mathrm{e}^{-rt}$. En un periodo la amplitud disminuye en un factor $\mathrm{e}^{-0.25 \times
4} = \mathrm{e}^{-1} = 0.37$. (d) Un forzamiento: las fluctuaciones aleatorias de los alisios, que
sacuden el sistema una y otra vez, y la \omterm{def:b2:climate-variability:bjerknes}{retroalimentación de Bjerknes}, que hace el amortiguamiento más débil o
incluso negativo.
\end{solution}

\begin{solution}{exo:b2:climate-variability:11}
(a) $A = R^2 \,\Delta\lambda\,(\sin\ang{5} - \sin(\ang{-5})) = (6.371 \times 10^6)^2 \times 0.8727
\times 0.1743 = \qty{6.17e12}{m^2}$. (b) $Q = 1025 \times 3992 \times 100 \times 2 \times 6.17
\times 10^{12} = \qty{5.05e21}{J}$. (c) Seis meses son $182.6 \times 86\,400 = \qty{1.578e7}{s}$;
por metro cuadrado, $8.18 \times 10^8/1.578 \times 10^7 = \qty{51.9}{W/m^2}$, mantenidos durante medio
año. (d) Del propio océano: unos alisios más débiles hacen subir menos agua fría en el este, y el agua cálida
del oeste, almacenada por encima de la termoclina más profunda, se extiende hacia el este; buena parte del calentamiento
puede proceder de calor que ya estaba en el océano, desplazado más que ganado.
\end{solution}

\begin{solution}{exo:b2:climate-variability:12}
(a) El bucle devuelve $0.3g$, después $0.3g^2$, y así sucesivamente: el total es $0.3(1 + g + g^2 + \dots) =
0.3/(1 - g)$, una serie geométrica que converge para $g < 1$: \qty{0.6}{K}, \qty{1.5}{K} y
\qty{3.0}{K} para $g = 0.5$, 0.8 y 0.9. (b) La serie diverge: la anomalía crecería sin
límite; el crecimiento lo deben detener procesos que el bucle deja fuera, el amortiguamiento por el flujo de calor
en superficie y el vaciado de la reserva de agua cálida que describe el oscilador de recarga. (c)
Cuanto más cálida es la superficie del mar, más fácilmente un pequeño calentamiento desencadena convección profunda y lluvia,
cuyo calentamiento impulsa los vientos; es la estación en que los efectos del ENOS son máximos.
\end{solution}

\begin{solution}{pb:b2:climate-variability:1}
\textbf{1.} $(0.26 + 0.75 + 1.13)/3 = \qty{0.713}{\celsius}$.
\textbf{2.} MJJ $(0.75 + 1.13 + 1.52)/3 = 1.133$; JJA $(1.13 + 1.52 + 1.80)/3 = 1.483$; JAS $(1.52 +
1.80 + 2.06)/3 = 1.793$.
\textbf{3.} ASO $(1.80 + 2.06 + 2.26)/3 = 2.040$; SON $(2.06 + 2.26 + 2.42)/3 = 2.247$; OND $(2.26
+ 2.42 + 2.36)/3 = 2.347$; NDE $(2.42 + 2.36 + 2.33)/3 = 2.370$.
\textbf{4.} Redondeados a un decimal, 0.7, 1.1, 1.5, 1.8, 2.0, 2.2, 2.3 y 2.4: la tabla.
\textbf{5.} Un solo mes lleva el ruido del tiempo; el ENOS dura muchos meses, y la media de tres
elimina buena parte del ruido conservando la señal lenta.
\textbf{6.} En JJA 1997, con 1.483, redondeado a 1.5; culminó en 2.4 en NDE.
\textbf{7.} De ASO 2018 a MJJ 2019, y de SON 2019 a FMA 2020.
\textbf{8.} Diez trimestres, y seis.
\textbf{9.} En ASO 2020; dentro de 2020 solo cuatro trimestres, de ASO a NDE, están en $-0.5$ o por debajo, uno
menos de lo que exige la regla, y DEF 2021 la completa.
\textbf{10.} Un El~Niño, de SON 2018 a MAM 2019, siete trimestres (AMJ 2019 baja a 0.37); el
episodio de La~Niña empieza en AMJ 2020 ($-0.55$), cuatro trimestres antes de donde la sitúa el ONI, en ASO.
\textbf{11.} No: el índice relativo nunca supera 0.29 a finales de 2019. Todo el océano tropical estaba
entonces más cálido que su periodo base en aproximadamente medio grado; la región Niño 3.4 estaba cálida en
términos absolutos, pero apenas más que los trópicos que la rodean.
\textbf{12.} Once.
\textbf{13.} 8, 7, 10, 4, 5, 6, 12, 5, 9 y 3 años.
\textbf{14.} El más corto es de 3 años (de 2023 a 2026), el más largo de 12 (de 1997 a 2009).
\textbf{15.} $69/10 = 6.9$ años de ambas maneras: la suma de los intervalos sucesivos es el tiempo desde el
primer inicio hasta el último, $2026 - 1957 = 69$ años.
\textbf{16.} Ya ha durado cinco trimestres, de MAM a JAS, en \qty{0.5}{\celsius} o por encima, y
ha culminado en 2.2, por encima de 1.5; su inicio, MAM 2026, es el primer trimestre de la sucesión, y los trimestres posteriores
solo pueden prolongarla, no desplazar su comienzo.
\textbf{17.} No, y no tiene por qué serlo: los 25 episodios de todas las intensidades llegan aproximadamente cada tres años, como dicen las preguntas frecuentes;
los fuertes son un subconjunto, así que son más raros, y sus intervalos van de 3 a 12 años.
\textbf{18.} $\tau = 1025 \times 3992 \times 50/16.1 = \qty{1.27e7}{s} = \qty{147}{\day}$.
\textbf{19.} Capa de mezcla: $1025 \times 3992 \times 50 = \qty{2.05e8}{J.m^{-2}.K^{-1}}$; columna de aire:
$(10^5/9.81) \times 1004.5 = 10\,194 \times 1004.5 = \qty{1.02e7}{J.m^{-2}.K^{-1}}$, 20 veces
menos.
\textbf{20.} $1.024 \times 10^7/(1025 \times 3992) = \qty{2.5}{m}$ de agua de mar.
\textbf{21.} $1.024 \times 10^7/16.1 = \qty{6.36e5}{s} = 7.4$~días.
\textbf{22.} Una anomalía llevada solo por el aire se amortigua en una semana aproximadamente, una de la capa de mezcla en
unos cinco meses, y la capa cálida por encima de la termoclina es aún más profunda: solo el océano retiene
el calor el tiempo suficiente para llevar una anomalía de una estación a la siguiente, mientras que la atmósfera aporta
los vientos que acoplan a ambos.
\textbf{23.} Desde 1950, los episodios fuertes de El~Niño han empezado en promedio cada \textbf{6.9 años},
aproximadamente siete.
\end{solution}
